\chapter{评测、短板与后续工作}

\section{哪些数字能说，哪些不能说}

\begin{keynote}{在阅读本章数字之前}
下列结果的性质是\emph{系统自检}，证明的是“实现与设计一致”，
而不是“输出内容正确”。原因已在 3.2 节说明：
知识库当前核实率仅百分之十九。
\end{keynote}

以五十组用例、百分之三十注入率运行的评测结果见表 \ref{tab:metrics}。

\begin{table}[htbp]\centering\small
\caption{三项指标与消融对照}\label{tab:metrics}
\begin{tabular}{@{}p{2.8cm}p{2.2cm}p{3cm}p{2.2cm}@{}}
\toprule
配置 & 幻觉率 & 适配规则一致性 & 覆盖率 \\
\midrule
完整流程 & \emph{0.00\%} & 100\% & 100\% \\
关闭审核 & 8.87\% & 100\% & 100\% \\
关闭辩论 & 0.00\% & 100\% & 100\% \\
榜题目标 & 小于 5\% & 不低于 85\% & 不低于 90\% \\
\bottomrule
\end{tabular}
\end{table}

消融对照是评测中最值钱的一栏。只报告“幻觉率为零”没有说服力，
报告“关掉审核是 8.87\%，开着是 0\%”才是能站得住的因果论证。

\begin{keynote}{“适配”那一项是自证的，别把它当效果指标}
系统按预先登记的规则计算资源难度，这一项又用同一条规则去校验，
所以它必然接近百分之百。它证明的是“实现与规则相符”，
证明不了“规则本身合理”。

因此在所有对外材料中，这一项一律写作\emph{“适配规则一致性”}，
不写“适配准确率”。答辩时若说成“我们的适配准确率达到百分之百”，
评委追问一句“你拿什么当真值”就会塌。
评测脚本在这一项达到百分之百时会主动打印警告，
就是为了防止有人顺手把它抄进汇报材料。
\end{keynote}

所有阈值集中在配置文件中。
\emph{修改配置文件要在版本库里留一次独立提交，提交时间必须早于评测批次的时间戳。}
规则先定后测，而不是测完再凑定义，守住了才叫科学定标。

\section{只有人能供给的两样东西}

三项指标当中适配那一项是自证的，幻觉率也只能证明“与知识库一致”。
\emph{任何自动化检验都绕不开这层循环，因为判据和被判对象同源。}
打破循环只有一个办法：引入一个与系统无关的判断源，也就是人。

系统为此提供两件仪器，它们不产生判断，只把样本抽好、把干扰去掉、
把一致性算出来。

\subsection{知识库导入器}

输入为 Markdown，一条切片一个二级标题。
选择 Markdown 而非表格格式，是因为正文里逗号、引号、换行都很常见，
表格转义一旦出错就是静默截断。

\emph{校验在入库时做，不在事后做。} 一条不合规的切片一旦进了库，
就会被检索到、被生成引用、被审核认证为“有依据”，
再想揪出来得翻整条链路；入口拦住只需要一次拒绝。
批内只要有一条被拒，整批不入库——半批入库会让编号与核实率都对不上账。

拒收项包括：伪造出处、知识点不存在、长度越界、编号重复、状态自相矛盾。
其中\emph{伪造出处}最为要紧——一个像真的假出处会让所有人误以为内容有据可查，
它没有任何视觉破绽，比明显编造更危险。

\begin{keynote}{踩坑：规范和语料对不上}
写这个工具时发现，规范初稿写“切片 150 到 300 字”，
而\emph{骨架版二十六条切片全部落在 71 到 104 字，没有一条符合我们自己写的规范}。

改了阈值而不是改切片。理由是“一条切片只讲一件事”这条更根本——
一个工业事实用中文讲清楚通常就是八九十字，硬凑到 150 字只能靠注水，
而注水的句子会拉低审核的证据覆盖率，\emph{反过来伤害幻觉检测}。
\end{keynote}

\subsection{人工评分工具}

盲评有三条设计。\emph{隐藏系统结论}：评分表里不出现系统给的难度值、
掌握概率、共识等级——看到系统答案再打分，打出来的是对系统的认同度，
不是独立判断，有测试专门检查评分表不泄漏这些字段。
\emph{打散顺序}：每位评分者拿到的顺序不同，由固定种子决定，可复现。
\emph{低一致不取平均}：一致性系数低于 $0.6$ 时，
工具会明确写出“请两位评分者一起过一遍分歧条目，对齐标准后重标”——
对不齐的两组数取平均，只会得到一个看着精确的假数。

\begin{keynote}{机器造仪器，人出判断}
这与核实智能体是同一个思路。机器能做的是把样本抽好、把证据摆到人面前、
把一致性算准，最后那一下必须由人来点。
凡是让机器代替人下最终判断的设计，都会在某个环节退化成“模型说它查过了”。
\end{keynote}

\section{已知短板清单}

答辩时主动讲短板，比被评委问出来强得多。本节建议原样进入方案文档。

\begin{table}[htbp]\centering\small
\caption{已知短板与补法}\label{tab:gaps}
\begin{tabular}{@{}p{2.4cm}p{3.4cm}p{6.2cm}@{}}
\toprule
短板 & 现状 & 影响与补法 \\
\midrule
知识库未核实 & 核实率百分之十九 & \emph{压倒其余全部短板}。知识库错则审核闸把错误认证为正确，所有幻觉率数字都只是自检 \\
H5、H6 类幻觉 & 检出率仅 67\% & 纯规则看不见语义层面的过度泛化与似真外推，需增加蕴含判断层 \\
知识库规模 & 仅二十六条切片 & 实际需要三百条以上，覆盖率指标全靠它 \\
题库规模与质量 & 四十七道，十九道待改 & 自适应要发挥作用需每知识点五到六道 \\
适配指标自证 & 结构性达到 100\% & 需要一项\emph{独立于生成规则}的检验，或人工评分 \\
共识分层倍数 & 离线桩产出 & 部分由桩的构造决定，真模型下必须重测 \\
知识追踪参数 & 文献常用初值 & 需用真实作答数据重新拟合 \\
生成题难度 & 标签是“请求值” & 无作答数据可标定，需积累数据后做项目分析 \\
辩论对齐 & 字符二元组相似度 & 认不出改写，应换用句向量或蕴含判断 \\
会话持久化 & 存于内存 & 进程重启即丢失，演示够用，不是生产实现 \\
\bottomrule
\end{tabular}
\end{table}

\section{接入大模型}

复制配置模板并填写本机密钥，随后运行体检脚本：

\begin{code}
Copy-Item .env.example .env
# 编辑 .env，分别填写 MiniMax 与 DeepSeek 的 API Key
# MiniMax 国内站: https://api.minimaxi.com/v1（国际站: https://api.minimax.io/v1）
# DeepSeek: https://api.deepseek.com
# 默认模型 MiniMax-M3，强任务模型 deepseek-v4-pro

python3 -m evalkit.doctor
\end{code}

\emph{未配置密钥时系统照样运行}，自动退回离线桩——
接入模型是增强项而非必需项，演示现场断网也不会挂掉。

体检脚本包含六项检查，其中最要紧的是\emph{命题实测}。
之所以需要它，是因为接入模型最常见的三类故障——
JSON 模式不被支持、模型名写错、触发限流——
在全流程中的表现\emph{都是“生成结果为空”}，极难定位。

系统固定使用两个模型：命题等高风险任务优先调用
\texttt{deepseek-v4-pro}，其余任务使用 \texttt{MiniMax-M3}。
强模型在重试后仍不可用时自动降级到 MiniMax；MiniMax 再失败时，
上层才回退固定题库与规则版。统计结果会记录实际模型与降级次数。

冷启动时，路由只按任务类型选择模型；每个模型积累五个样本后才允许动态重排。
备选模型的成功率高出至少二十个百分点，或成功率接近而当前模型平均延迟达到
备选模型两倍时，路由才切换优先级。连续三次可熔断故障会打开断路器，六十秒后
仅放行一次 half-open 探测；成功则恢复，失败则继续熔断。两家供应商凭据独立时，
某一家的 401 不会重试该模型，但允许降级到另一家；统一网关共用凭据时才直接失败。

比赛演示完全不需要付费。免费档真正的约束不是“能不能用”，
而是每分钟请求数与总额度，代码为此做了主动限速、调用次数硬上限、
请求缓存三条保护。其中调用次数上限尤其重要：
额度悄悄耗尽后系统会一路回退规则版，
\emph{“能跑但已降级”是最坏的一种故障}，因为使用者并不知道自己在演示哪个版本。
每一次 HTTP 尝试（包括失败与重试）都消耗 \texttt{AGENTEDU\_BUDGET\_CALLS}，
配置上限时必须为失败留出预算。

服务启动后，操作者可以在本机通过
\texttt{GET /api/model-status} 查看两个模型的健康、尝试、成功率、平均延迟和
自动降级次数。该接口是只读且脱敏的：不返回 API Key、Authorization 请求头、
提示词、端点 URL、请求体或模型响应。单元测试只使用本地假端点和测试值，
绝不读取真实密钥。

\begin{keynote}{踩坑：额度用完把整条链路炸掉}
命题环节原先没有捕获模型异常，调用失败的异常一路冒到编排层，整个演示白屏。
而这个场景下正确的行为是\emph{降级}：命题本来就是增强项，
没有它系统照样能用固定题库跑完。
\emph{演示现场额度用完还能继续演，和当场白屏，是两回事。}
\end{keynote}

\section{排期与分工}

按九月五日提交倒推的排期见表 \ref{tab:sched}。

\begin{table}[htbp]\centering\small
\caption{四周排期}\label{tab:sched}
\begin{tabular}{@{}p{1.8cm}p{10cm}@{}}
\toprule
周次 & 任务 \\
\midrule
第 1 周 & 验通演示服务；用摄入工具录入真实资料；扩充知识点表；改完待办的十九道题 \\
第 2 周 & 接入真模型并运行体检；重跑全部评测；补前端资源正文渲染 \\
第 3 周 & 人工标注与一致性计算；三张图表打磨；单元测试补到三百个以上 \\
第 4 周 & 撰写测评报告与方案文档；制作十分钟视频；打包提交 \\
\bottomrule
\end{tabular}
\end{table}

知识库建设是纯体力活，但不能交给不懂领域的人。
建议由一人负责确定切片边界与知识点归属，此人必须懂领域；
两人负责录入；确定边界的那位再\emph{逐条复核}知识点归属——
归属错了，检索的硬过滤就失效了。
\emph{第一周就要启动，拖到第三周一定来不及。}

演示视频须在十分钟以内，\emph{制作时间要留足三天而不是三小时}，
写好脚本再录，即兴录制一定会超时。

最后是两条合规提醒。项目中的三份学习者画像\emph{全部是合成数据}，
提交材料中应明确写出，零成本即可对应数据合规条款。
竞赛总则要求作品拥有完整独立知识产权，
而\emph{知识库是最容易出问题的地方}——
教材与手册的原文大段照抄存在版权风险，
建议以自行改写的技术要点为主，来源字段标明依据的资料。
这件事要早处理，不要等到打包时才想起来。

\bibliographystyle{gbt7714-2005}
\bibliography{refs}
